\section{An all-order criterion for the first lens-shell interval}\label{app:energy}
Let the action on the two complex coordinates have rotation weights \((1,r)\) modulo \(p\), \(\gcd(r,p)=1\). For each \(t=0,1\) obeying \(1+r+tp=0\bmod2\), put \(a=(1+r+tp)/2,\ b=(1-r+tp)/2\). Then \(m_+(0)=2\mathbf1_{b=0\bmod p}\) and \(m_-(0)=2\mathbf1_{a=0\bmod p}\). For \(p\ge2\) at most one side has a first-shell doublet. Changing \(r\) to \(-r\) exchanges \(a,b\) and reverses the eta sign. In this generator convention \(L(7,2)\) has \((a,b)=(5,3)\), \(\eta_0=-4/7\), the mirror of the row with \(+4/7\) in Appendix~\ref{app:spheres}.

When \(\delta_0\ne0\), the finite group acts in one \(SU(2)\) factor only. The character formula reduces to an elementary sum whose sign is that of \(\delta_0\):
\[
 \frac{\eta_0}{\delta_0}=\frac1p\sum_{j=1}^{p-1}
                 \frac1{2-2\cos(2\pi j/p)}=\frac{p^2-1}{12p}.
\]
The last identity follows from logarithmic derivatives of \((z^p-1)/(z-1)\) at \(z=1\). Setting \eqref{eq:energy} to zero gives \(X_*=3/[2(1-\eta_0/\delta_0)]\); an interior zero requires \(0<\eta_0/\delta_0<2/5\). Exactly \(p=2,3,4\) satisfy it. When \(\delta_0=0\), \(D=2X\eta_0\) is identically zero in this interval or has no zero there. For \(L(7,2)\) the next-shell imbalance gives the first positive zero at \(7/2\); both mirror conventions have that position. The finite register contains six lifts with isolated first-interval zeros, ten identically zero there and 418 without an interior zero.
